% Reconstructed from the compiled lecture-note PDF.
\chapter{The matrix language: SL2(R) without prior Lie theory}
The matrix language: without SL2(R) prior Lie theory


\section{Area-preserving matrices}

\begin{displaytext}
Let \
SL2(R) = \{g ∈M2(R) : det g = 1\} . \
This is a group under matrix multiplication. Geometrically, g ∈SL2(R) is an orientation-preserving \
linear transformation that preserves signed area:
\end{displaytext}

\begin{displaytext}
det(gu, gv) = det(g) det(u, v) = det(u, v).
\end{displaytext}

The corresponding space of infinitesimal matrices is


\begin{displaytext}
( a b ! ) \
sl2(R) = \{X ∈M2(R) : tr X = 0\} = : a, b, c ∈R . \
c −a
\end{displaytext}

The terminology “Lie group” and “Lie algebra” will be explained in Appendix A; for now these are concrete matrix sets.


The matrix 0 −1! J = 1 0


generates rotations because cos t −sin t ! exp(tJ) = . sin t cos t


\begin{displaytext}
We write \
π \
R = exp J \
3 \
for counterclockwise rotation by 60◦.
\end{displaytext}

\section{Every multipoint is an SL2(R)-image of the standard one}

\begin{displaytext}
Let (sj) and (σj) be two normalized multipoints. Since s0 and s2 are linearly independent, there is \
a unique linear map g satisfying \
gs0 = σ0, gs2 = σ2.
\end{displaytext}

\begin{displaytext}
The sum relations express every other sj as a linear combination of s0, s2, and the same relations \
hold for σj. Therefore \
gsj = σj
\end{displaytext}

\begin{displaytext}
for all j. Finally, \
det(gs0, gs2) det(σ0, σ2) \
det(g) = = = 1. \
det(s0, s2) det(s0, s2) \
Thus g ∈SL2(R).
\end{displaytext}

Applying this at every time gives the basic representation.


\begin{statementbox}{Theorem}
Theorem 4.1 (Matrix representation of a multicurve). For a normalized multicurve there is a unique curve g : I →SL2(R) such that
\end{statementbox}

\begin{displaytext}
σj(t) = g(t)s∗j for all j, t.
\end{displaytext}

The differentiability of g is the same as that of the multicurve.


Six synchronized curves have now been replaced by one curve in a three-dimensional matrix manifold.


\section{Body velocity}

\begin{displaytext}
Differentiate σj(t) = g(t)s∗j: \
σ′j(t) = g′(t)s∗j.
\end{displaytext}

Since g(t) is invertible, define


\begin{displaytext}
X(t) = g(t)−1g′(t), equivalently g′(t) = g(t)X(t).
\end{displaytext}

This X is the velocity measured in the moving frame determined by g. It belongs to sl2(R).


To prove tracelessness, differentiate det g(t) = 1. Jacobi’s formula gives


\begin{displaytext}
d \
0 = det g = det(g) tr(g−1g′) = tr X. \
dt
\end{displaytext}

Thus a(t) b(t) ! X(t) = . c(t) −a(t)


\begin{displaytext}
The boundary tangents are \
σ′j(t) = g(t)X(t)s∗j. \
Equivalently, because σj = gs∗j, \
σ′j = gXg−1σj. \
The conjugate gXg−1 is the same infinitesimal motion written in the fixed laboratory coordinates.
\end{displaytext}

Background interlude


The map X 7→gXg−1 is called the adjoint action and is denoted Adg X. Nothing abstract is hidden here: it is ordinary matrix conjugation. The commutator


[X, Y ] = XY −Y X


is the corresponding infinitesimal operation.


\section{Translating the star condition into inequalities}

\begin{displaytext}
The standard points are \
cos(jπ/3)! \
s∗j = . sin(jπ/3)
\end{displaytext}

\begin{displaytext}
Because multiplication by g preserves open cones, the star condition for σ′j = gXs∗j is equivalent to \
a condition on Xs∗j before applying g. \
For each j, the vector Xs∗j must lie in the cone spanned by s∗j+1 and s∗j+2. Solving the corresponding \
linear systems yields three linear forms \
√ \
c + 3a \
ρ0(X) = √ , (4.1) \
3 \
√ \
c − 3a \
ρ1(X) = √ , (4.2) \
3 \
+ c ρ2(X) = −3b√ . (4.3) \
2 3
\end{displaytext}

The six cone conditions reduce, by central symmetry, to


\begin{displaytext}
ρ0(X) > 0, ρ1(X) > 0, ρ2(X) > 0.
\end{displaytext}

These are the matrix form of the star inequalities.


\begin{displaytext}
They imply several useful signs: \
√ \
c > 0, 3 |a| < c, 3b + c < 0, tr(JX) = b −c < 0.
\end{displaytext}

The final inequality will guarantee that the area integrand is positive.


\section{A determinant identity}

For a b ! X = , c −a


we have det X = −a2 −bc.


A direct calculation using Equations (4.1) to (4.3) gives


\begin{displaytext}
det X = ρ0ρ1 + ρ1ρ2 + ρ2ρ0.
\end{displaytext}

Since all ρj are positive, det X > 0.


This positivity permits a convenient reparametrization.


\section{Normalizing time by det X = 1}

If we replace the parameter t by an increasing parameter τ, then


\begin{displaytext}
dg dt \
= gX . \
dτ dτ \
Thus X is multiplied by the positive scalar dt/ dτ, and its determinant is multiplied by the square \
of that scalar. Choose \
dt 1 \
= dτ p det X(t). \
Then the reparametrized matrix, again denoted X, satisfies
\end{displaytext}

det X(t) = 1.


This is analogous to arclength parametrization, but it is adapted to the matrix geometry rather than Euclidean speed.


For a traceless 2 × 2 matrix, the characteristic polynomial is


λ2 + det X.


\begin{displaytext}
Under det X = 1, Cayley–Hamilton gives \
X2 = −I, X−1 = −X.
\end{displaytext}

This simple identity is used repeatedly.


\section{The trace form}

\begin{displaytext}
Define \
⟨X, Y ⟩= tr(XY ), X, Y ∈sl2(R).
\end{displaytext}

This symmetric bilinear form is nondegenerate. It has the invariance property D gXg−1, gY g−1E = ⟨X, Y ⟩,


\begin{displaytext}
and the commutator identity \
⟨X, [Y, Z]⟩= ⟨[X, Y ], Z⟩. \
For X ∈sl2(R), \
⟨X, X⟩= tr(X2) = −2 det X.
\end{displaytext}

Hence the determinant-one surface is the quadratic surface


⟨X, X⟩= −2.


The sign is not a mistake: the trace form is indefinite, not an ordinary Euclidean inner product.


\section{Reconstructing the critical hexagon from X}

The directions of the six support edges at the standard multipoint are Xs∗j. The intersection of the lines through s∗j and s∗j+1 in these directions gives a vertex of the current critical hexagon. Solving the two linear equations yields


\begin{displaytext}
ρj+2(X) pj = s∗j + Xs∗j = s∗j+1 −ρj(X) Xs∗j+1. det X det X
\end{displaytext}

When det X = 1, this simplifies further. Thus X does not merely encode velocity; together with the current multipoint it reconstructs the support hexagon.


\begin{insightbox}{Geometric intuition}
\end{insightbox}

The three-dimensional matrix X simultaneously records:


• the six tangent directions of the synchronized boundary curves; • the shape of the current critical hexagon; • the position of the system in a hyperbolic state space introduced later.


The star inequalities select one connected region of the determinant-one quadratic surface.


\section{Endpoint conditions}

In circle representation, g(0) = I.


\begin{displaytext}
After one sixth of the boundary has been traversed, the jth curve must arrive where the (j + 1)st \
curve started. Therefore \
g(tf)s∗j = s∗j+1 = Rs∗j,
\end{displaytext}

so g(tf) = R.


Smooth matching of tangents gives


X(tf) = R−1X(0)R.


Extending the solution by this symmetry produces a periodic orbit:


\begin{displaytext}
X(t + 3tf) = X(t), g(t + 6tf) = g(t).
\end{displaytext}

The periodicity of a minimizer will be decisive at the end of the proof.


\begin{insightbox}{Chapter summary}
\end{insightbox}

\begin{displaytext}
The six-point geometry is encoded by one area-preserving matrix g(t), and its body velocity \
X = g−1g′ is traceless. Convexity becomes three linear star inequalities ρj(X) > 0. These \
imply det X > 0, so time can be normalized by det X = 1. The endpoint conditions are \
rotational, making every admissible minimizer periodic after extension.
\end{displaytext}
